% Generated draft: requires mathematical and source review.
\documentclass[11pt]{article}
\usepackage[margin=25mm]{geometry}
\usepackage{fontspec}
\setmainfont{Latin Modern Roman}
\usepackage{amsmath,amssymb,mathtools,unicode-math}
\setmathfont{Latin Modern Math}
\usepackage{enumitem,needspace,xurl,hyperref}
\setlength{\parindent}{0pt}
\setlength{\parskip}{5pt}
\setlength{\emergencystretch}{3em}
\newcommand{\N}{\mathbb{N}}
\newcommand{\Z}{\mathbb{Z}}
\newcommand{\Q}{\mathbb{Q}}
\newcommand{\R}{\mathbb{R}}
\newcommand{\C}{\mathbb{C}}
\newcommand{\F}{\mathbb{F}}
\newcommand{\A}{\mathbb{A}}
\newcommand{\PP}{\mathbb{P}}
\newcommand{\E}{\mathbb{E}}
\newcommand{\eps}{\varepsilon}
\newcommand{\dd}{\,\mathrm{d}}
\begin{document}
\section*{UnsolvedMath ID: 20003105}
\textbf{Source problem code:} catalogue entry
\subsection{Definitions and mathematical statement}
A. Papadopoulos's source question asks whether the length-spectrum metric on Teichm\"uller space is Finsler and, if so, for its infinitesimal norm. The surface class, normalization of the metric, and smoothness required of ``Finsler'' are not specified. The supplied reduction uses the usual finite-type hyperbolic Teichm\"uller setting and a possibly nonsmooth continuous norm.

Fix such a surface \(S\), and write \(\ell_X(\lambda)\) for hyperbolic length of a nonzero measured lamination. In the normalization used by the aid, put
\[
 d_{\mathrm{Th}}(X,Y)=\log\sup_{\lambda\ne0}\frac{\ell_Y(\lambda)}{\ell_X(\lambda)},\qquad
 d_{\mathrm{ls}}(X,Y)=\max\{d_{\mathrm{Th}}(X,Y),d_{\mathrm{Th}}(Y,X)\}.
\]
A convention inserting an overall factor \(1/2\) rescales all corresponding norms and path distances; it must not be mixed with these formulas. Thurston's asymmetric norm is
\[
 F_X(v)=\sup_{[\lambda]}\dd_X\log\ell_\lambda(v).
\]
The aid states that the metric derivative of \(d_{\mathrm{ls}}\) forces the reversible norm
\[
 N_X(v)=\max\{F_X(v),F_X(-v)\}
       =\sup_{[\lambda]}|\dd_X\log\ell_\lambda(v)|.
\]
For a differentiable path \(\gamma\) with \(\gamma(0)=X\), \(\dot\gamma(0)=v\), this is the stated limit of \(d_{\mathrm{ls}}(X,\gamma(t))/|t|\).

Let
\[
 D_N(X,Y)=\inf_{\gamma:X\to Y}\int_0^1N_{\gamma(t)}(\dot\gamma(t))\,\dd t,
\]
with the appropriate piecewise smooth or absolutely continuous admissible paths. The reported general inequality is \(D_N(X,Y)\geq d_{\mathrm{ls}}(X,Y)\). A local norm formula alone does not establish equality of these global distances.

For the aid's minimizing-path criterion, suppose the \(N\)-distance between the pair is attained by an admissible path. Put \(a=d_{\mathrm{Th}}(X,Y)\), \(b=d_{\mathrm{Th}}(Y,X)\). It reports:
\begin{itemize}
\item If \(a>b\), equality \(D_N(X,Y)=d_{\mathrm{ls}}(X,Y)\) is equivalent to a path \(\gamma:X\to Y\) which is a forward Thurston geodesic up to reparametrization and satisfies \(F_{\gamma(t)}(\dot\gamma(t))\geq F_{\gamma(t)}(-\dot\gamma(t))\) almost everywhere.
\item If \(b>a\), the analogous condition uses a reverse Thurston geodesic, with \(F(-\dot\gamma)\geq F(\dot\gamma)\) almost everywhere.
\item If \(a=b\), equality requires one path geodesic in both directions, with \(F(\dot\gamma)=F(-\dot\gamma)\) almost everywhere. Such a compatible bi-geodesic is also sufficient.
\end{itemize}
The directional geodesic condition means its corresponding asymmetric length equals the endpoint distance; a parametrization with unit speed is not an extra requirement. A compatible path is automatically an \(N\)-minimizer. Conversely, deduction of a single compatible path from equality uses attainment of the infimum; without attainment the formulation must use approximating paths instead. The record does not spell out the existence/properness argument for minimizers or the full surface hypotheses.

In the intended continuous, possibly nonsmooth Finsler sense, the forced local norm leaves the global question whether \(D_N=d_{\mathrm{ls}}\) for every pair. Under the required minimizer-existence hypotheses this becomes existence of the compatible paths above for every pair. The aid describes that global existence question as still open; it does not establish that the length-spectrum metric is Finsler.
\subsection{Short English statement}
The supplied reduction determines the only possible reversible infinitesimal norm, but its intrinsic path metric may exceed the endpoint length-spectrum distance. Equality requires compatible directional geodesics; tied forward and reverse distances require a single equal-speed bi-geodesic. The global existence issue is left open.
\subsection{Sources}
\begin{itemize}
\item Local UnsolvedMath registry, permanent ID \(20003105\), exact source code \texttt{catalogue entry}; canonical \texttt{aim-topology-notes.json}, \texttt{notes[192]}, tag \texttt{problem}.
\item AIM workshop ``Lipschitz metric on Teichmueller space,'' identified in the supplied verification as the 2012 workshop, Problem 3.2, attributed to A. Papadopoulos: \url{https://aimath.org/pastworkshops/lipschitzteichproblems.pdf}. The record calls the canonical question exact. This supplied PDF URL was not fetched.
\item The original literature field is empty. The record supplies no complete authoritative bibliographic reference for the infinitesimal metric-derivative formula, minimizer existence, or global open-status assessment.
\item Research attempt 1: result \texttt{reduction}, problem status at run \texttt{open}, rigor label \texttt{proved}, novelty confidence \texttt{low}; classification \texttt{PARTIAL-PROGRESS}, registry status \texttt{partially\_solved}, suggested difficulty L4. Created and updated 14 August 2026. These are supplied dated assessments, not acceptance of the original Finsler claim.
\item Attribution: Ulam AI / UnsolvedMath Contributors, CC BY 4.0. The machine-generated research aid is not peer reviewed and requires independent verification before citation.
\end{itemize}
\end{document}
